% chapters/ch06-integrity-constraints.tex
\chapter{Integrity Constraints and Correctness}

\begin{learningobjectives}
In this chapter you will learn how constraints make a relational design reliable.
You will be able to classify constraints (entity, referential, domain, business), decide which rules can be enforced by the database, and document constraints as part of the schema specification.
\end{learningobjectives}

\section{Constraints are part of the meaning}
A schema without constraints is only a storage layout.
It can hold values, but it does not protect meaning.
Relational databases become powerful because they are not passive containers: they can reject invalid states.

When a system relies only on application logic to maintain correctness, the database becomes an unsafe shared memory.
Multiple services, scripts, and manual interventions will eventually bypass the intended rules.
Constraints exist to prevent this drift.

\begin{definitionbox}
A \term{constraint} is a condition that every valid database state must satisfy.
Constraints exist to prevent invalid states and preserve domain meaning over time.
\end{definitionbox}

In this chapter we continue to avoid SQL.
We focus on the conceptual role of constraints and how to specify them clearly, independent of implementation.

\section{Constraint types with examples}
Not all constraints are the same.
Classifying them helps you know what the database can enforce directly and what may require additional mechanisms.

\begin{figure}[htbp]
\centering
\includegraphics[width=0.75\textwidth]{constraint-types.pdf}
\caption{Four constraint categories clarify what can be enforced directly and what must be specified explicitly.}
\end{figure}

\subsection{Entity integrity}
Entity integrity is the requirement that identity is real.
Primary keys must uniquely identify rows and must not be missing.
This is not optional: without it, relationships cannot be trusted.

Even in a prose specification, you should state this clearly.
For each relation, identity constraints include:
the primary key attributes, whether they are mandatory, and whether they are stable.

\subsection{Referential integrity}
Referential integrity is the requirement that references point to real things.
If \rel{Appointment} references \rel{Patient}, then every appointment must point to an existing patient.
Otherwise you are modelling relationships to non-existent entities, which breaks the semantics of the model.

Referential integrity also forces you to decide lifecycle behaviour.
When a referenced entity is removed, what happens to referencing rows?
Sometimes deletion is forbidden; sometimes references cascade; sometimes a record is archived instead of deleted.
Even if you postpone the exact policy, you should recognise that it is part of correctness.

\subsection{Domain constraints}
Domain constraints restrict allowed values.
They can be as simple as “age is non-negative”, or as meaningful as “due date must be on or after checkout date”.
They often capture basic validity that is easy to enforce and easy to forget.

Domain constraints are valuable because they prevent silent corruption.
If invalid values are allowed in, the system will later contain “impossible” states that break analysis and reporting.

\subsection{Business constraints}
Business constraints express domain-specific rules that go beyond identity, references, and simple value checks.
These are the most important constraints from a correctness perspective, and also the most often missed.

A clinic rule like “a doctor cannot have two appointments at the same time” is a business constraint.
A library rule like “a copy cannot have two active loans at the same time” is also a business constraint.
These constraints define what the system \emph{means}, not merely what it stores.

\begin{examplebox}
\textbf{University constraint mix}\par\smallskip
In a course registration model:
\begin{itemize}[leftmargin=*,nosep]
  \item Entity: Student.student\_id is present and unique.
  \item Referential: Enrollment.student\_id references Student.
  \item Domain: Enrollment.grade is one of the allowed grade symbols.
  \item Business: A student cannot enroll in the same course twice in the same term.
\end{itemize}
\end{examplebox}

\section{Feasibility of enforcement}
A practical discipline is to treat each business rule as a constraint candidate.
For each rule, you ask three questions.

First: can the database enforce it directly as a static constraint?
Some rules map cleanly to uniqueness, not-null requirements, and foreign keys.

Second: if it cannot be enforced as a simple constraint, can it be enforced with a stronger mechanism?
Some rules require temporal reasoning or overlap checks.

Third: if enforcement cannot realistically live in the database, can violations be detected reliably?
A rule is still useful if you can detect when it is violated and react.

This discipline makes the design defensible.
It also makes responsibilities explicit.

\begin{principlebox}
\textbf{Rule \(\rightarrow\) enforcement discipline}\\[2pt]
For each business rule:
\begin{itemize}[leftmargin=*,nosep]
  \item Can the database enforce it directly?
  \item If not, can it enforce it with a stronger mechanism?
  \item If not, how do we detect violations reliably?
\end{itemize}
\end{principlebox}

\begin{principlebox}
\textbf{Constraints are part of the schema}\\
If a rule is not stated as a constraint, it is not part of the design.
Write the rules before you choose where to enforce them.
\end{principlebox}

There is a temptation to say “the database enforces everything”.
That is not always realistic.
Some constraints involve complex temporal logic, external systems, or probabilistic rules.
Others are policy decisions that evolve.

The right goal is not maximal enforcement at the database level.
The right goal is: enforce what you can, and explicitly handle what you cannot.
A rule that is unenforced and undocumented is the worst case.
A rule that is enforced in the application but clearly specified and monitored can be acceptable.

\begin{notebox}
A practical rule: enforce identity and referential integrity in the database by default.
For domain and business constraints, enforce in the database when feasible and document the rest with detection strategies.
\end{notebox}

\section{Constraint catalogs}
Catalogs make constraints visible.
They also provide a checklist for validation later.

\subsection{Clinic constraints}
\begin{itemize}[leftmargin=*,nosep]
  \item C1: Patient.patient\_id is mandatory and unique.
  \item C2: Doctor.doctor\_id is mandatory and unique.
  \item C3: Appointment.appointment\_id is mandatory and unique.
  \item C4: Appointment.patient\_id references an existing Patient (mandatory).
  \item C5: Appointment.doctor\_id references an existing Doctor (mandatory).
  \item C6: Appointment.start\_time is mandatory.
  \item C7: Appointment.duration\_min is mandatory and positive.
  \item C8: A doctor cannot have overlapping appointments.
  \item C9: A patient cannot have overlapping appointments with the same doctor.
  \item C10: Appointment.room\_id references an existing Room when present.
  \item C11: A room cannot host overlapping appointments.
  \item C12: Appointment.status is one of \{scheduled, completed, cancelled\}.
\end{itemize}

\subsection{Library constraints}
\begin{itemize}[leftmargin=*,nosep]
  \item L1: Member.member\_id is mandatory and unique.
  \item L2: Title.title\_id is mandatory and unique.
  \item L3: Copy.copy\_id is mandatory and unique.
  \item L4: Copy.title\_id references an existing Title (mandatory).
  \item L5: Copy.barcode is unique when present.
  \item L6: Loan.loan\_id is mandatory and unique.
  \item L7: Loan.member\_id references an existing Member (mandatory).
  \item L8: Loan.copy\_id references an existing Copy (mandatory).
  \item L9: Loan.due\_date must be on or after Loan.checkout\_date.
  \item L10: A copy has at most one active loan at any time.
  \item L11: Reservation.member\_id references an existing Member (mandatory).
  \item L12: Reservations for a title are ordered by request\_date.
\end{itemize}

\section{Constraint writing style guide}
Constraints should be written as testable sentences.
A good constraint is precise, scoped, and observable.

Write constraints in the form: “For every X, condition Y holds.”
Avoid vague language like “usually” or “should”.

Even without SQL, you should specify constraints as part of the schema.
A reader should be able to answer: “what does this model allow?”

A compact documentation format works well:
\begin{itemize}[leftmargin=*,nosep]
  \item list relations and keys,
  \item list foreign key references,
  \item list uniqueness constraints,
  \item list domain and business constraints as numbered statements.
\end{itemize}

This style is close to how real systems are specified in design docs.
It also makes your later SQL chapter easier to write: SQL becomes the translation of an already-clear specification.

\begin{examplebox}
\textbf{Constraint block example}\par\smallskip
For \rel{Copy}:
\begin{itemize}[leftmargin=*,nosep]
  \item C1: Copy.copy\_id is mandatory and unique.
  \item C2: Copy.title\_id references an existing Title (mandatory).
  \item C3: Copy.barcode is unique when present.
  \item C4: Copy.status is one of \{available, on\_loan, retired\}.
\end{itemize}
\end{examplebox}

\section{Common mistakes}
Constraint mistakes are usually mistakes of silence.

Mistake: Failing to write the rule at all.\\
Fix: Record each rule as a numbered, testable constraint.

Mistake: Treating a business rule as a convenience.\\
Fix: Decide whether it is truly a rule; if yes, enforce or detect it.

Mistake: Hiding constraints in unrelated attributes.\\
Fix: State the constraint explicitly (e.g., uniqueness, no-overlap).

Mistake: Mixing domain and business constraints in prose only.\\
Fix: Separate value rules (domain) from relationship rules (business).

\section{Mini checklist}
Before you sign off a schema, check:
\begin{itemize}[leftmargin=*,nosep]
  \item Are identity and reference rules explicit for every relation?
  \item Are domain rules stated for key attributes with limited values?
  \item Are business rules written as constraints, not prose hints?
  \item Is each constraint placed where it can be enforced or detected?
  \item Are lifecycle rules (delete, archive) clearly described?
  \item Do constraint catalogs cover the critical entities and events?
  \item Are “active” and “overlap” defined unambiguously?
\end{itemize}

\begin{chaptersummary}
Constraints are not optional decorations; they are part of the meaning of a relational design.
Entity and referential integrity provide the foundation for identity and relationships.
Domain constraints prevent basic invalid values.
Business constraints preserve domain-specific correctness, such as preventing double booking or concurrent active loans.
A good design makes enforcement explicit, either directly in the database or through documented detection and control mechanisms.
\end{chaptersummary}

\begin{exercisebox}
Answer in full sentences.
\begin{enumerate}[label=\arabic*.]
  \item For a domain you know, write 8 business rules and classify each as entity, referential, domain, or business constraint.
  \item Pick 3 of those rules and explain whether the database can enforce them directly, and why.
  \item In a library model, propose a precise definition of “active loan” and rewrite the “one active loan per copy” rule using that definition.
  \item In a clinic model, explain why interval-based appointments make “no double booking” harder than slot-based appointments.
  \item Write a constraint specification block (C1, C2, ...) for one relation in your schema.
  \item For a university model, define a business constraint that uses term or semester and state how it could be enforced.
  \item Choose one rule that cannot be fully enforced by the database and describe how you would detect violations.
  \item Write a domain constraint for a status field and explain what values are allowed.
  \item Propose a constraint that prevents overlapping room bookings and state the needed attributes.
\end{enumerate}
\end{exercisebox}

In the next chapter we validate designs without SQL.
The constraint catalogs you wrote here will become inputs to validation.
